% Consolidated earlier incentive results; insert before reproducible evidence.
\section{Finite-pool mechanisms and limits of restricted rules}\label{hist:appendix}
This appendix collects nonredundant results from the earlier incentive
campaigns. All guarantees use the fixed-owner, binding-underreport,
payment-only model of Section~\ref{sec:model}. The subsidy factor is
$\beta=B-1$, not $B$. The mechanisms and their restricted-class comparisons
remain useful even where later theorems improve the asymptotic frontier.
The weighted-glove model, ordinary Shapley value, and continuous Banzhaf
functional are established ingredients; the comparisons below make no
historical-priority claim.

\subsection{Aggregate prices and an exact restricted-class tradeoff}
\begin{proposition}\label{hist:aggregate}
For $T>0$, write $C_j=c(N_j)$ and define
\[
 p_j=\frac1{C_j}\int_0^T\prod_{\ell\ne j}(1-t/C_\ell)\,dt,
 \qquad y_i=c_ip_j\quad(i\in N_j).
\]
The rule covers coalitions, is own-report monotone, protects $H_i/k$, and
has total payment at most $(k+1)T/2$. This budget bound is a supremum for
this rule. For two pools and $\alpha\in[0,1]$, the optimal worst-case
subsidy factor among rules $y_i=c_ip(C_1,C_2)$ in pool one and
$y_i=c_iq(C_1,C_2)$ in pool two is
$\max\{1/2,\alpha\}$, provided the same nonnegative prices are used for
\emph{every} finite ownership partition and the rule protects $\alpha H$.
\end{proposition}
\begin{proof}
Differentiating $P(t)=\prod_j(1-t/C_j)$ and integrating gives
$\sum_jp_j=P(0)-P(T)=1$. Thus payments to a coalition are a convex
combination of its pool totals. For an owner reporting $x$ with fixed
same-pool remainder $a$, set
$g(t)=\prod_{\ell\ne j}(1-t/C_\ell)_+$. Its payment is
\[
 \frac{x}{x+a}\int_0^{x+a}g(t)\,dt.
\]
Both factors are nonnegative and nondecreasing in $x$. Assign zero at
$x=a=0$; if another pool has zero total all payments stay zero. This proves
monotonicity at every boundary as well.
The complementary-coalition constraint gives
$H_i\le T-v_c(N\setminus\{i\})\le c_iT/C_j$.
For the last inequality, it is equality as an upper bound when $C_j=T$;
otherwise the marginal is $(c_i-C_j+T)_+\le c_iT/C_j$.
Also $p_j\ge T/(kC_j)$, since every factor in the integral is at least
$1-t/T$. Choose a pool of total $T$. Its payment is at most $T$;
every other pool receives at most $\int_0^T(1-t/T)dt=T/2$.
Totals $(1,M,\ldots,M)$ approach this bound as $M\to\infty$.

For the two-pool upper bound, the displayed rule gives $\beta=1/2$ and
$H/2$. The auxiliary rule $z_i=c_iT/C_j$ gives $H$, covers coalitions,
and spends $2T$. Its own payment is $x\min\{1,D/(x+a)\}$, where $D$ is
the other pool total, so it is nondecreasing. For $\alpha\ge1/2$, the
convex combination $(2-2\alpha)y+(2\alpha-1)z$ gives $\alpha H$ and
$\beta\le\alpha$; for smaller $\alpha$, use $y$ itself.

For the lower bound, one owner per pool at totals $(1,1)$ gives
$p,q\ge\alpha$ and $p+q\ge1$. One price is at least
$r=\max\{1/2,\alpha\}$; suppose it is $p$.
Raising that owner's report to $M$ gives $Mp(M,1)\ge r$.
Now partition the left total into owners of capacities $1,M-1$.
Partition-independent prices are unchanged. The unit left owner together
with the right owner imply $p(M,1)+q(M,1)\ge1$. Hence total payment is
at least $1+(M-1)p(M,1)\ge1+r(1-1/M)$. Let $M\to\infty$.
This argument uses the stated comparison \emph{across} ownership
partitions; it is not an unrestricted fixed-owner lower bound.
\end{proof}

For clarity, these pool totals are themselves the established absolute
continuous Banzhaf functional \cite[Definition~6.11]{Kurz}. For
$F(z)=\min_j C_jz_j/T$, the zero endpoint is zero and the other endpoint
has tail integral
\[
 T\operatorname{BZI}_j(F)
 =\E\min\{C_j,\min_{\ell\ne j}C_\ell U_\ell\}
 =\int_0^T\prod_{\ell\ne j}(1-t/C_\ell)dt=C_jp_j.
\]
Here the coordinates are \emph{pools}. The owner-coordinate identification
in Section~\ref{sec:banzhaf} concerns a different continuous game.

\subsection{What ordinary Shapley scaling can and cannot achieve}
Let $\phi$ be the ordinary individual Shapley value, with fixed weights
$|U|!(n-|U|-1)!/n!$ on predecessor sets $U\subseteq N\setminus\{i\}$.
Nonnegativity and efficiency follow by monotonicity and telescoping along
each arrival order. Each own marginal is nondecreasing in its owner's
capacity, so $\phi_i$, and every nonnegative constant multiple of it, is
own-report monotone; this truthfulness argument is established in
\cite[Proposition~1]{Huang}.

\begin{proposition}\label{hist:shapley-half}
In every finite two-pool game,
$\phi_i\ge M_i/2\ge H_i/2$ and $\phi(S)\ge v_c(S)/2$.
Both constants $1/2$ are sharp over finite ownerships. Consequently a
constant multiple $\lambda\phi$ covers every coalition at every profile
and ownership if and only if $\lambda\ge2$. Even a profile-dependent
scalar multiplier cannot have a uniform budget factor less than two.
\end{proposition}
\begin{proof}
Fix an owner of capacity $a$, same-pool remainder $A$, and opposite total
$D+A$. Put $f_a(z)=\min\{a,z_+\}$. A predecessor subset with opposite-minus-
same capacity $z$ gives marginal $f_a(z)$, while its complement gives
$f_a(D-z)$. For all real $z,D$,
\[
 f_a(z)+f_a(D-z)\ge f_a(D).
\]
If $D\le0$ this is immediate; if either argument is at least $a$ it is
also immediate. Otherwise the positive parts are untruncated and their
sum is at least $D$. Complementary predecessor sets have equal Shapley
weights, so $2\phi_i\ge f_a(D)=M_i$. The profile $(1\mid1)$ attains
$\phi_i=H_i/2$. If $L$ is a globally minimum-total pool, every $i\in L$
has $M_i=c_i$; hence
$\phi(S)\ge\phi(S\cap L)\ge c(S\cap L)/2\ge v_c(S)/2$.

For sharpness of the coalition constant, choose an integer $q\ge2$,
$m=q^4$, $r=q^4-q^3$, and a game with $m$ left unit owners, $r$ right
unit owners, and one right owner $b$ of capacity $m$. Let $S$ contain all
left owners and $b$, so $v_c(S)=T=m$. Independent uniform arrival times
represent a uniform random order. For a left unit arriving at time $t$,
its marginal is one if $b$ has arrived; otherwise it is one precisely
when $Y>X$, with independent $X\sim\operatorname{Bin}(m-1,t)$ and
$Y\sim\operatorname{Bin}(r,t)$. Thus
\[
 \phi_{\rm left}=\tfrac12+\eta,
 \qquad \eta=\int_0^1(1-t)\Prob(Y>X)dt.
\]
Writing $D_0=m-1-r=q^3-1$, the difference has mean $-D_0t$ and variance
at most $m/2$. Chebyshev gives
$\Prob(Y>X)\le\min\{1,m/(2D_0^2t^2)\}$, whose integral is at most
$\sqrt{2m}/D_0$. At $b$'s arrival, its marginal is $(X-Y)_+$ for counts
with parameters $m,r$; Cauchy--Schwarz bounds its expectation by
$(m-r)t+\sqrt{m/2}$. Therefore
\[
 \frac{\phi(S)}m\le\frac12+
 \frac{\sqrt2q^2}{q^3-1}+\frac1{2q}+\frac1{\sqrt2q^2}\longrightarrow\frac12.
\]
This proves an infimum over finite games, without splitting any actual
owner. Coverage forces any scalar multiplier to be at least
$v_c(S)/\phi(S)$, and efficiency identifies that multiplier with its
budget factor, proving the last two claims.
\end{proof}

\begin{proposition}\label{hist:four-shapley}
With exactly two owners in each of two pools, $(6/5)\phi$ covers all
coalitions, is own-report monotone, protects $3H/5$, and spends $6T/5$.
The multiplier $6/5$ is sharp for constant Shapley scaling on this domain.
It fails for some five-owner games.
\end{proposition}
\begin{proof}
Write capacities $(a,b\mid c,d)$, set $m_{ij}=\min(c_i,c_j)$, and
$w_i=v_c(N\setminus\{i\})$. Directly grouping the predecessor subsets gives
\[
 \phi_1=\frac{m_{13}+m_{14}+w_4-m_{23}+w_3-m_{24}+w_2}{12}
          +\frac{T-w_1}{4},
\]
with the other coordinates obtained by pool-preserving relabelings.
Only cross pairs and triples need a nontrivial coverage check.
For $S=\{2,3\}$,
\[
 6\phi(S)=m_{23}+w_4+w_1-m_{14}+3T-w_2-w_3.
\]
Exchange pools if needed so $a+b\le c+d$. Then $T=a+b,w_1=b,w_2=a$,
and
\[
 3T+w_1+w_4-w_2-w_3-m_{14}
 \ge3b+\min(a+b,c)\ge4\min(b,c).
\]
Here $w_3\le a+b$ and $m_{14}\le a$. Thus $\phi(S)\ge5v_c(S)/6$.
For $S=\{1,2,3\}$ the desired inequality, after using efficiency, is
\[
 K=9T-7w_4-w_1-w_2-w_3-m_{14}-m_{24}+m_{13}+m_{23}\ge0.
\]
Put $L=a+b,R=c+d$. If $L\le R$ and $c\ge L$, this expression is
$[L-\min(L,d)]+[a-\min(a,d)]+[b-\min(b,d)]\ge0$.
If $L\le R$ and $c<L$, bounding $w_3,m_{14},m_{24}$ by $L,a,b$ gives
$K\ge6(L-c)+\min(a,c)+\min(b,c)-c\ge0$.
If $L>R$, use $w_4=c,w_3=d$,
$\min(x,c+d)\le\min(x,c)+d$ for $x=a,b$, and $m_{14},m_{24}\le d$;
then $K\ge2c+4d\ge0$. These cases include zeros and ties.
Every other cross pair or triple follows by relabeling.
Proposition~\ref{hist:shapley-half} supplies the claimed fairness floor.
At $(2,1\mid1,1)$,
$\phi=(7/12,1/4\mid7/12,7/12)$, so a pair containing the unit left
owner and a right owner has worth one and payment $5/6$.
At $(1,1\mid1,1,2)$ the Shapley vector is
$(41/60,41/60\mid11/60,11/60,4/15)$, obtained by the same predecessor
formula. The two left owners and the capacity-two right owner have worth
two but receive only $49/25$ after scaling by $6/5$.
\end{proof}

\subsection{The ownership-dependent exponential two-pool rule}
\begin{theorem}\label{hist:exponential}
For arbitrary finite owners in two pools, let $d_i=\min(c_i,T)$.
Call a minimum-total pool short, put
$\rho=d(N_{\rm long})/T$, and $p=\tfrac12e^{1-\rho}$.
Pay $(1-p)c_i$ to short owners and $pd_i$ to long owners, with zero
payments at $T=0$. This rule covers coalitions, is own-report monotone,
protects $M_i/2\ge H_i/2$, and has subsidy at most $1/(2e)$.
The constant is attained for this rule. Every admissible rule, even
without a floor, has worst-case subsidy at least $1/8$ on two pools
with two fixed owners each.
\end{theorem}
\begin{proof}
Capping preserves worth by Lemma~\ref{lem:cap}, and $\rho\ge1$, so the
payments are a convex combination of the two capped pool totals. At a
tie, each owner is paid $c_i/2$, independently of the selected short pool.
Writing $u=\rho-1\ge0$, total payment is
$T[1+ue^{-u}/2]\le T[1+1/(2e)]$, with equality at $(1,1\mid1)$.
Short owners have $M_i=c_i$ and $1-p\ge1/2$. For a long owner, let $a$
be the total of its same-pool companions. If $a\ge T$ its marginal is
zero. Otherwise its companions are uncapped, and with $m=d_i$,
$M_i=T-a\le m\le T$ and $u=(m-M_i)/T$. Since $t e^{-t/T}$ increases
on $[0,T]$, its payment is at least $M_i/2$.

For own-monotonicity, fix a left owner reporting $x$, same-pool remainder
$a$, and other pool total $D>0$. In the long region $x+a\ge D$, put
$m=\min(x,D)$ and $\gamma=\sum_{h\ne i,\,h\in L}\min(c_h,D)$.
Its payment is $\tfrac12e^{1-\gamma/D}m e^{-m/D}$, nondecreasing in $x$.
In the short region $0<x+a\le D$, put $A=x+a$ and
$\rho(x)=\sum_{h\in R}\min(c_h/A,1)$.
Away from finitely many clipping points,
$0\ge\rho'\ge-\rho/A$, and the derivative of $x(1-p)$ is
\[
 1-p+xp\rho'\ge1-p(1+\rho)\ge0,
\]
where the last inequality follows because
$(1+\rho)e^{1-\rho}/2$ decreases from one on $[1,\infty)$.
The payment is continuous at all clipping points and at $A=D$, where
both formulas give $x/2$. If $D=0$ payments remain zero; if $A=0$ the
assigned zero payment is no greater than any later payment.

For the lower bound, use fixed owners $0,1\mid2,3$ and profiles
$c=(1,1\mid1,1)$, $a=(1,1\mid1,2)$, $b=(2,1\mid1,1)$.
At $a$, coverage of $\{0,2\},\{1,2\},\{0,1,3\}$ implies
$2y(a)(N)-y_3(a)\ge4$. The budget $y(a)(N)\le2+2\beta$ gives
$y_3(a)\le4\beta$. At $b$, coverage of
$\{1,2\},\{1,3\},\{0,2,3\}$ similarly gives $y_0(b)\le4\beta$.
Own-monotonicity and the base coalition $\{0,3\}$ give
$1\le y_0(c)+y_3(c)\le y_0(b)+y_3(a)\le8\beta$.
\end{proof}
The resulting interval $[1/8,1/(2e)]$ holds for the two-pool minimax
subsidy with any fixed floor fraction $0\le\alpha\le1/2$.
Neither endpoint is asserted to be the unrestricted optimum. The price
here depends on individual clipping, so Proposition~\ref{hist:aggregate}'s
partition-independent-price lower bound does not apply. The smaller
budget does not dominate Proposition~\ref{hist:four-shapley}'s stronger
$3H/5$ fairness guarantee.

\subsection{Small-stage guarantees and failed exponential extensions}
\begin{proposition}\label{hist:softmax}
For $k=2$ or $3$, put $\rho_j=\sum_{i\in N_j}d_i/T$, $u_j=\rho_j-1$,
$w_j=e^{-u_j}$, $p_j=w_j/\sum_\ell w_\ell$, and $y_i=d_ip_j$.
This rule is own-report monotone, covers coalitions, protects $M_i/k$,
and has exact worst-case subsidy $b_k$, where
$b_ke^{b_k}=(k-1)/e$. The same unit-rate formula fails own-report
monotonicity when $k=4$.
\end{proposition}
\begin{proof}
Coverage is immediate from capping and $\sum p_j=1$.
For an owner with $M_i>0$, same-pool remainder $a<T$, and $m=d_i$, one
has $u_j=(m-M_i)/T$ and $\sum w_j\le k$. Hence
$y_i\ge m e^{-(m-M_i)/T}/k\ge M_i/k$ by the preceding monotonicity of
$t e^{-t/T}$. The assertion is immediate if $M_i=0$.
In the long region of an own report, $T$ is fixed, and differentiation
with respect to $m$ gives $p_j[1-(m/T)(1-p_j)]\ge0$.
In the short region, $T=x+a$, $u_j=0$, and
$0\ge\rho_\ell'\ge-\rho_\ell/T$ off clipping points. Therefore
\[
 \frac{d}{dx}(xp_j)
 \ge p_j\left[1-\sum_{\ell\ne j}p_\ell\rho_\ell\right]\ge0,
\]
since the last bracket is nonnegative precisely when
$\sum_{\ell\ne j}u_\ell e^{-u_\ell}\le1$, and each term is at most
$1/e$. For $k\le3$, $(k-1)/e<1$. Continuity handles clipping points
and the bottleneck seam; the zero boundary is handled by $0\le y_i\le x$.

The subsidy is $F=\sum_j u_je^{-u_j}/\sum_j e^{-u_j}$, with some
$u_s=0$. For $b=b_k$, the elementary maximum
$(u-b)e^{-u}\le e^{-1-b}$ implies
\[
 \sum_j(u_j-b)e^{-u_j}\le-b+(k-1)e^{-1-b}=0.
\]
Thus $F\le b$. Equality holds when $u_s=0$ and every other $u_j=1+b$:
use one unit owner in pool $s$ and three owners of capacity $(2+b)/3<1$
in each other pool. Here $b<1$ follows from $k\le3$.

For failure at $k=4$, use one owner reporting $x$ in the first pool and
three owners of capacity $2/3$ in each other pool. Near $x=1$ the first
owner's payment is
$y(x)=x/[1+3e^{1-2/x}]$, whose derivative at one is
$(1-3/e)/(1+3/e)^2<0$. Thus some strict underreport increases payment.
\end{proof}
A different extension also fails: assigning each nonshort pool price
$e^{-u_j}/k$ and the residual to the short pool is not truthful even for
three pools. At $(x\mid1/2,1/2\mid1/2,1/2)$, the first payment is
$1/3$ when $x=1$, but at $x=99/100$ it is
\[
 x\left[1-\tfrac23 e^{-(1-x)/x}\right]
 \ge x(1-2x/3)=1683/5000>1/3,
\]
using $e^{-q}\le1/(1+q)$. These are obstructions to the stated formulas,
not to all mechanisms with more pools.

\begin{proposition}\label{hist:small-stage-lower}
For $k\ge3$ and two owners per pool, every admissible rule has
\[
 \beta\ge\max\left\{\frac{k-1}{2(k+3)},\frac{k-2}{2k+3}\right\}.
\]
In particular the lower bounds are $1/6$ at $k=3$ and $3/14$ at $k=4$;
the latter exceeds the two-pool upper bound $1/(2e)$.
\end{proposition}
\begin{proof}
Average over pool permutations and within-pool swaps, preserving every
guarantee on the same fixed ownership. The first displayed bound is
Theorem~\ref{thm:unrestricted} with $r=2,t=2$.
For the second, take profile $X$ with one pool $(1,1)$ and $k-1$ pools
$(1,2)$. Let $a$ be the former owners' common payment, and $b,d$ the
unit and capacity-two payments in each mixed pool. Coalitions selecting
all unit owners or both short owners and every capacity-two owner give
\[
 a+(k-1)b\ge1,\quad 2a+(k-1)d\ge2,\quad
 2a+(k-1)(b+d)\le2+2\beta.
\]
Subtracting the first inequality from the budget and comparing with the
second gives $a\ge1-2\beta$. At the profile $Z$ of all mixed pools,
write the payments $b_Z,d_Z$. Then $kb_Z\ge1$ and
$k(b_Z+d_Z)\le3+3\beta$, so $kd_Z\le2+3\beta$.
Raising one unit owner in $X$'s short pool to two produces $Z$;
own-monotonicity gives $d_Z\ge a$. Thus
$k(1-2\beta)\le2+3\beta$. Finally $e>5/2>7/3$ proves
$3/14>1/(2e)$ without numerical approximation.
\end{proof}

\subsection{An obstruction to capped-total-only prices}
\begin{proposition}\label{hist:capped-totals}
Suppose payments have the form $y_i=d_ip_j(\rho_1,\ldots,\rho_k)$,
where $\rho_j=d(N_j)/T$ and the nonnegative prices do not depend on
within-pool capacity distributions. Fixed owner counts and pool labels
may enter the prices. Coalition coverage and own-monotonicity imply,
for $k=q^2$ and $n$ owners per pool with $q\mid n$, the budget bound
\[
 B\ge\frac{k}{2\sqrt{k}-1}>\frac{\sqrt{k}}2.
\]
Prices need not sum to one. Zero reports are allowed.
\end{proposition}
\begin{proof}
At the profile with every owner reporting $1/n$, all ratios equal one
and the budget gives $\sum_jp_j\le B$. Choose a pool $j$ with
$p_j\le B/k$. Within this same fixed ownership, concentrate its capacity
in one owner, leaving its others zero and every other pool unchanged.
At the concentrated report one, the ratio vector is still all ones.
Lowering this owner's report to $x=1/q$ therefore gives
$xp_j(\rho^{\rm low})\le B/k$. A coalition containing this owner and
$n/q$ owners in each other pool has worth $x$ and payment
$x\sum_\ell p_\ell(\rho^{\rm low})$, so that sum is at least one.
The low-profile total payment is consequently at least
\[
 xp_j+\sum_{\ell\ne j}p_\ell
 \ge1-(1-x)p_j\ge1-\frac{(1-x)B}{kx}.
\]
Its budget is at most $Bx$, proving
$B\ge[x+(1-x)/(kx)]^{-1}=k/(2\sqrt{k}-1)$.
No independently strategic clones or comparison across ownerships is
used here. The restriction to ratio-only prices is essential; the
owner-coordinate Banzhaf rule uses the full within-pool distribution.
\end{proof}
